\section{制度得失：接口越多，谁来承担与检验}

汉的制度收获，不在于拥有一部无所不能的中央机器，而在于多次把原本分散的资源接到共同秩序：关中仓储和郡县遗产支撑开国，削藩后中枢更能核验王国军政，选官和文书把地方评价送入官署，财政措施扩大转运和筹费能力，边塞与使节将远方信息带回朝廷。东汉又在战争之后重建了雒阳、簿籍和道路的连接。

这些收获从来没有免费。征发、货币调整、边郡军费、迁徙、度田和军队供给落在地区与家庭身上的重量并不相同；现存材料又最少保存承担者的连续声音。中枢为应急设置的辅政、内廷、外戚、宦官、边将和州郡军权，也会把本该服务公共协调的节点变成可继承、可交易的政治资源。

因此，汉的损失不宜概括为“中央太弱”或“地方太强”。王国、家族和地方中介曾提供兵员、信息与保护；问题在于它们的资源是否还能被放进共同、可检验的责任链。新朝的改革暴露命令无法自动穿过地方，东汉的宫廷冲突暴露检验接口的困难，汉末的区域竞争则显示同一批接口已经不能服从同一套优先次序。

\begin{turningpoint}[衡量制度，不以成败贴标签]
\term{国家能力}问的是朝廷能否让命令、粮食、人员和信息接续，而不是地图是否涂成同一种颜色。

\term{财政可持续性}问的是远征、赏赐、屯田和救济由谁支付，不能只计一场胜利或一项诏令。

\term{社会代价与政治检验}问的是谁能对征发和任命提出异议，谁在危机中失去保护，谁把公共节点变成自己的资源。用这些尺度连读西汉、新、东汉，才能看到制度的成就与脆弱性原本由同一组连接产生。
\end{turningpoint}
